%% kapa: exam=1 round=32 year=2021 period=1 form=A subject=economics subject_ko=경제학원론 seq=27 number=67
%% source: Q-Net 감정평가사 기출문제 (cst00309) artlSeq=5208397 file=extracted/1008/5208397/A형/2021년도 제32회 감정평가사 1차 1교시 A형.pdf
%% conversion: kordoc 4.24.51 -> tools/gen1.py (auto); review=proofread-sonnet; flags=none
\begin{question}{67}
소비자물가지수에 관한 설명으로 옳지 않은 것은?
\begin{choices}
  \item 기준연도에서 항상 100이다.
  \item 대체효과를 고려하지 못해 생계비 측정을 왜곡할 수 있다.
  \item 가격 변화 없이 품질이 개선될 경우, 생계비 측정을 왜곡할 수 있다.
  \item GDP 디플레이터보다 소비자들의 생계비를 더 왜곡한다.
  \item 소비자가 구매하는 대표적인 재화와 서비스에 대한 생계비용을 나타내는 지표이다.
\end{choices}
\end{question}
\answer{4}
